\documentclass[11pt,a4paper]{ctexart}

\usepackage{amsmath,amssymb,mathtools}
\usepackage{geometry}
\usepackage{booktabs}
\usepackage{hyperref}
\usepackage{xcolor}
\usepackage{enumitem}

\geometry{margin=2.4cm}
\hypersetup{
  colorlinks=true,
  linkcolor=blue!60!black,
  urlcolor=blue!60!black,
  citecolor=blue!60!black
}
\setlist[itemize]{leftmargin=2em}
\setlist[enumerate]{leftmargin=2em}

\title{Model, Planning 与 Synthetic Rollout}
\author{}
\date{\today}

\begin{document}
\maketitle

\section{总览}

在 world model 和 WAM 文献中，\textbf{model}、\textbf{planning} 和 \textbf{synthetic rollout} 是三个紧密相关但不等价的概念。

一句话区分：
\[
  \text{model} = \text{学到的环境近似},
\]
\[
  \text{planning} = \text{利用模型选择动作},
\]
\[
  \text{synthetic rollout} = \text{在模型内部生成的虚拟轨迹}.
\]

它们之间的关系是：先学习一个模型，再用该模型展开未来轨迹，最后基于这些未来轨迹进行动作选择或策略训练。

\section{Model：学到的世界近似}

在序列决策问题中，真实环境的动力学可以写为：
\[
  p_{\mathrm{env}}(s_{t+1}\mid s_t,a_t),
\]
其中 $s_t$ 是真实状态，$a_t$ 是动作。

但真实状态通常不可见，机器人只能看到观测 $o_t$，例如图像、深度图或 proprioception。因此实际学习的 world model 常写为：
\[
  p_\theta(o_{t+1}\mid o_{\le t},a_t),
\]
或多步形式：
\[
  p_\theta(o_{t+1:t+H}\mid o_{\le t},a_{t:t+H-1}).
\]

这里 $p_\theta$ 是用数据训练出的模型，用来近似真实环境：
\[
  p_\theta \approx p_{\mathrm{env}}.
\]
如果模型还预测奖励或价值，可以扩展为：
\[
  p_\theta(o_{t+1},r_t\mid o_{\le t},a_t),
\]
或
\[
  p_\theta(o_{t+1:t+H},r_{t:t+H}\mid o_{\le t},a_{t:t+H-1}).
\]

\subsection{模型的作用}

World model 的核心作用不是直接执行动作，而是提供一个可查询的近似环境。给定当前观测和候选动作，模型回答：
\[
  \text{如果执行这些动作，未来可能会怎样？}
\]

因此 model 是 planning 和 synthetic rollout 的基础。

\section{Synthetic Rollout：模型内部的虚拟轨迹}

真实 rollout 是智能体在真实环境中执行动作得到的轨迹：
\[
  \tau =
  (o_0,a_0,o_1,a_1,\ldots,o_T).
\]
其中每个下一步观测来自真实环境：
\[
  o_{t+1}\sim p_{\mathrm{env}}(o_{t+1}\mid o_{\le t},a_t).
\]

Synthetic rollout 则不是从真实环境采样，而是从学到的模型采样：
\[
  \hat{o}_{t+1}\sim p_\theta(o_{t+1}\mid o_{\le t},a_t).
\]
然后策略继续基于模型生成的观测选动作：
\[
  a_{t+1}\sim \pi(a_{t+1}\mid \hat{o}_{t+1}).
\]
再由模型预测下一步：
\[
  \hat{o}_{t+2}\sim p_\theta(o_{t+2}\mid o_{\le t},a_t,\hat{o}_{t+1},a_{t+1}).
\]
重复这个过程得到：
\[
  \hat{\tau}
  =
  (o_t,a_t,\hat{o}_{t+1},a_{t+1},\hat{o}_{t+2},\ldots,\hat{o}_{t+H}).
\]
这条轨迹就是 synthetic rollout，也可以叫 imagined rollout、model rollout 或 imagined trajectory。

\subsection{为什么叫 synthetic}

它是 synthetic 的原因是：轨迹中的未来观测不是传感器真实采集的，而是模型合成的：
\[
  \hat{o}_{t+k}\not\sim p_{\mathrm{env}},
  \qquad
  \hat{o}_{t+k}\sim p_\theta.
\]
因此 synthetic rollout 的质量完全依赖于模型是否接近真实环境。

\section{Planning：用模型选择动作}

Planning 的目标是选择一段动作，使未来累计回报最大。理想情况下，应该在真实环境动力学下求解：
\[
  a_{t:t+H-1}^*
  =
  \arg\max_{a_{t:t+H-1}}
  \mathbb{E}_{p_{\mathrm{env}}}
  \left[
    \sum_{k=0}^{H-1}\gamma^k r_{t+k}
  \right].
\]
但真实环境不可随意查询，机器人试错成本很高，所以用 world model 代替真实环境：
\[
  a_{t:t+H-1}^*
  =
  \arg\max_{a_{t:t+H-1}}
  \mathbb{E}_{p_\theta}
  \left[
    \sum_{k=0}^{H-1}\gamma^k \hat{r}_{t+k}
  \right].
\]

也就是说，planning 是在模型内部先试动作，再选择看起来最好的动作。

\subsection{采样式 planning}

很多 embodied AI 方法不直接解析求解最优动作，而是采样 $N$ 条候选动作序列：
\[
  a_{t:t+H-1}^{(i)}\sim q(a_{t:t+H-1}\mid o_{\le t},\ell),
  \qquad i=1,\ldots,N.
\]
对每条动作序列，用模型生成 synthetic rollout：
\[
  \hat{\tau}^{(i)}\sim p_\theta(\cdot\mid o_{\le t},a_{t:t+H-1}^{(i)}).
\]
再给每条 rollout 打分：
\[
  J^{(i)}
  =
  \sum_{k=0}^{H-1}\gamma^k \hat{r}_{t+k}^{(i)}
  \quad \text{或} \quad
  J^{(i)}=\hat{V}(\hat{o}_{t+H}^{(i)}).
\]
最后选择分数最高的动作序列：
\[
  i^*=\arg\max_i J^{(i)}.
\]
实际执行时通常只执行第一个动作或第一个 action chunk：
\[
  a_t = a_t^{(i^*)}.
\]
下一时刻重新观测环境并再次规划，这就是 receding horizon control 或 model predictive control 的思想。

\section{三者的关系}

\begin{center}
\begin{tabular}{lll}
\toprule
概念 & 数学对象 & 作用 \\
\midrule
Model &
$p_\theta(o_{t+1}\mid o_{\le t},a_t)$ &
近似环境动力学 \\
Synthetic rollout &
$\hat{\tau}\sim p_\theta$ &
在模型中生成虚拟未来 \\
Planning &
$\arg\max_a \mathbb{E}_{p_\theta}[\sum \gamma^k r]$ &
基于虚拟未来选择动作 \\
\bottomrule
\end{tabular}
\end{center}

因此：
\[
  \text{model 是工具，synthetic rollout 是模型产生的数据，planning 是使用这些数据做决策的过程。}
\]

\section{和 WAM 的关系}

在传统 world model 中，动作通常是条件：
\[
  p_\theta(o_{t+1:t+H}\mid o_{\le t},a_{t:t+H-1}).
\]
模型回答的是：给定动作之后，世界如何变化？

在 WAM 中，动作也可能是生成目标的一部分：
\[
  p_\theta(o_{t+1:t+H},a_{t:t+H-1}\mid o_{\le t},\ell).
\]
这意味着 WAM 不只是 rollout model，也是 action model。它可以直接生成：
\[
  (a_{t:t+H-1},\hat{o}_{t+1:t+H}).
\]

如果再加入价值预测，形式变为：
\[
  p_\theta(o_{t+1:t+H},a_{t:t+H-1},V_{t+H}\mid o_{\le t},\ell).
\]
这时 planning 可以通过 best-of-$N$ 采样完成：
\[
  (a^{(i)},\hat{o}^{(i)},\hat{V}^{(i)})\sim p_\theta,
  \qquad
  i^*=\arg\max_i \hat{V}^{(i)}.
\]
然后执行：
\[
  a_t=a_t^{(i^*)}.
\]

\section{Synthetic Rollout 的主要风险}

Synthetic rollout 的核心风险是模型误差：
\[
  p_\theta \neq p_{\mathrm{env}}.
\]
单步预测误差可能不大，但多步展开会产生误差累积：
\[
  \hat{o}_{t+1}\text{ 有误差}
  \Rightarrow
  \hat{o}_{t+2}\text{ 误差更大}
  \Rightarrow
  \hat{o}_{t+H}\text{ 可能偏离真实分布}.
\]
这通常称为 compounding error。

另一个风险是 model exploitation。规划器可能找到一条在模型中得分很高、但在真实世界中不可行的动作序列：
\[
  \arg\max_a \mathbb{E}_{p_\theta}[R]
  \neq
  \arg\max_a \mathbb{E}_{p_{\mathrm{env}}}[R].
\]
这也是为什么很多方法只做短 horizon rollout，或者不断用真实观测校正模型。

\section{一句话总结}

Model 是对真实环境的近似；synthetic rollout 是在这个近似环境中生成的虚拟未来；planning 是利用这些虚拟未来选择动作。

在 WAM 语境下，这三者进一步融合：模型不仅想象世界，还生成动作，并可能用价值函数对不同未来进行排序。

\end{document}
